% ===================================================================
%  TABEL Nr. 2 — Inimkeha. --- Vahetund. --- Mängud.
%  Traduction 2026 d'apres texto/io, controlee sur texto/fr et
%  texto/fi. Consigne : voir texto/et/10-tabel-01.tex.
%
%  LA PREMIERE MOITIE DE CETTE PAGE EST LE SEUL ENDROIT DE
%  L'ESTONIEN OU HUIT SIECLES DE SEIGNEURS ALLEMANDS N'ONT RIEN
%  LAISSE.
%
%  Pres d'un tiers du vocabulaire estonien vient du bas-allemand : le
%  tableau 1 l'a montre sur le mobilier de la classe. Le corps, lui,
%  est intact et finno-ougrien de bout en bout — « pea », « kere »,
%  « aju », « nägu », « silm », « kõrv », « suu », « keel »,
%  « kael », « rind », « süda », « selg », « käsi », « jalg »,
%  « veri », « nahk », « luu », « kand ». Pas un emprunt sur cette
%  liste. On emprunte le nom d'un objet qu'un marchand apporte ; on
%  n'emprunte pas le nom de sa propre main.
%
%  LA SECONDE MOITIE, LES JEUX, DIT AUTRE CHOSE, ET C'EST LA REPONSE
%  A LA COLONNE ROMANCHE.
%
%  Le romanche du tableau 2 avait pris « velo » a l'allemand de
%  Suisse et « chegls » et « stelzas » au meme voisin : la chose
%  arrivait du nord, le mot est arrive avec elle. L'estonien recoit
%  exactement les memes objets, par la meme route et a la meme date,
%  et n'en prend AUCUN mot :
%    la bicyclette est « jalgratas », la roue-a-pied ;
%    le football est « jalgpall », la balle-au-pied ;
%    la gymnastique est « võimlemine », de « võim », la force ;
%    l'escrime est « vehklemine », de « vehkima », agiter ;
%    la nage est « ujumine ».
%
%  LE PLANIFICATEUR NE CHANGE PAS LA ROUTE DE LA CHOSE, IL CHANGE
%  SEULEMENT SI LE MOT L'ACCOMPAGNE. C'est la precision que les
%  seize colonnes precedentes ne pouvaient pas donner : elles ne
%  connaissaient que des langues qui se sont faites toutes seules. Une
%  langue qu'on a decidee prefere le compose transparent a l'emprunt
%  opaque, parce qu'un compose S'ENSEIGNE — et le tableau 1 a dit qui
%  l'avait decide, et quand.
%
%  Le romanche avait fait un compose lui aussi pour le football,
%  « ballape ». Un tel accord entre deux langues sans rapport ne
%  s'explique pas par la planification : il s'explique par la chose.
%  Un ballon qu'on frappe du pied SE DECRIT, et deux langues qui le
%  decrivent tombent sur le meme compose. La regle vaut donc pour ce
%  qui ne se decrit pas, et la bicyclette en est le cas net.
% ===================================================================

%%K t02-apar-1 apar
\VUsaut{4.0mm}
\VUtitre{15.0pt}{60}{TABEL Nr. 2}
\VUsaut{2.0mm}
\VUtitre{11.4pt}{0}{Inimkeha. --- Vahetund.}
\VUtitre{11.4pt}{0}{Mängud.}

%%K t02-tit-0 sub
\VUtitre{13.2pt}{0}{Inimkeha.}
\VUsaut{2.4mm}
\VUcentre{10.2pt}{0}{\textit{(Lihtne loodusteaduse tund.)}}
\VUsaut{3.0mm}

%%K t02-tit-1 sub pos=t02-01-1
\VUpk{10.2pt}{40}{Pea.}
\VUsaut{3.0mm}

%%K t02-01-1 p
1. --- Inimkeha peamised osad on: \VUgras{pea}, \VUgras{kere} ja \VUgras{jäsemed}.

%%K t02-02-1 p
2. --- Pea koosneb kahest osast: \VUgras{koljust}, mis sisaldab \VUgras{aju} ja mida
katavad \VUgras{juuksed}, ning \VUgras{näost}.

%%K t02-03-1 p
3. --- Meie joonisel ei näe me koljut ega aju; kuid me näeme väga selgesti näo
tähtsaid osi.

%%K t02-04-1 p
4. --- Kõigepealt on siin \VUgras{laup}, mis annab märku tugevast mõistusest ja alati
tegusast mõttest. \VUgras{Meelekohad} on väga vabad. \VUgras{Silm} on elav ja laialt
avatud. Tihe \VUgras{kulm} ja pikad \VUgras{ripsmed} kaitsevad seda.

%%K t02-05-1 p
5. --- Siin on veel \VUgras{nina} ja \VUgras{sõõrmed}. Nina abil me nuusutame ja
hingame. Kummalgi pool nina on \VUgras{põsed}, kaugemal \VUgras{kõrvad}. Näo alumine
osa koosneb \VUgras{suust} ja \VUgras{lõuast}.

%%K t02-06-1 p
6. --- Me ei näe neid väga hästi, sest \VUgras{vuntsid} ja \VUgras{põskhabe} varjavad
neid meie eest. Kuid me teame, et suus on kaks \VUgras{huult}, \VUgras{ülalõug} ja
\VUgras{alalõug} koos oma \VUgras{hammastega}, ning \VUgras{keel}. Suu abil me räägime
ja sööme. Keele abil me maitseme toitu.

%%K t02-07-1 p
7. --- Silm, kõrv, nina ja suu on, nagu ka sõrmed, viie meele elundid:
\VUgras{nägemismeel}, \VUgras{kuulmismeel}, \VUgras{haistmismeel},
\VUgras{maitsmismeel}, \VUgras{kompimismeel}.

%%K t02-08-1 p
8. --- Pea on seega üks inimkeha huvitavamaid osi.

%%K t02-tit-2 sub
\VUsaut{4.4mm}
\VUpk{10.2pt}{40}{Kere.}
\VUsaut{3.0mm}

%%K t02-09-1 p
9. --- \VUgras{Kael} ühendab pea kerega. See hõlmab \VUgras{kurgu} ja \VUgras{kukla}.

%%K t02-10-1 p
10. --- Kere sisaldab ka palju eluliselt tähtsaid elundeid. \VUgras{Rinnas} on
\VUgras{kopsud}, millega me hingame, ja \VUgras{süda}, mis on elu kese. Kui süda
lakkab tuksumast, siis elusolend sureb.

%%K t02-11-1 p
11. --- \VUgras{Roided} hoiavad rinda.

%%K t02-12-1 p
12. --- Kere ülejäänud osad on \VUgras{külg}, \VUgras{selg}, \VUgras{õlad} ja
\VUgras{kõht}. Me heidame magama küljele või seljale, me kanname koormaid oma õlal.

%%K t02-13-1 p
13. --- Kõhus on \VUgras{magu} ja \VUgras{sooled}. Need aitavad meil toitu seedida.

%%K t02-tit-3 sub
\VUsaut{4.4mm}
\VUpk{10.2pt}{40}{Jäsemed.}
\VUsaut{3.0mm}

%%K t02-14-1 p
14. --- Meil on neli \VUgras{jäset}: kaks \VUgras{kätt} ja kaks \VUgras{jalga}. Kätega
me võtame, hoiame, tõstame ja korjame esemeid; jalgadega me seisame, kõnnime ja
jookseme.

%%K t02-15-1 p
15. --- \VUgras{Õlg} ühendab käe kehaga. Väga tugev \VUgras{lihas}, kakspealihas,
liigutab kätt, mille \VUgras{küünarnukk} ühendab \VUgras{käsivarrega}. \VUgras{Ranne}
moodustab \VUgras{käelaba} liigese, mis lõpeb \VUgras{sõrmede} ja \VUgras{küüntega}.
Käel eristame \VUgras{veeni} (ja arterit), milles voolab \VUgras{veri}.

%%K t02-16-1 p
16. --- Jala osad on: \VUgras{reis}, \VUgras{põlv} ja \VUgras{põlveõnnal},
\VUgras{sääremari}, mille \VUgras{liha} katab \VUgras{nahk}, \VUgras{pahkluu} ja
\VUgras{jalalaba}. Jala \VUgras{luud} on väga jämedad ja väga tugevad. Jalalaba otsas
on \VUgras{varbad}. Ülejäänud osad on \VUgras{tald} ja \VUgras{kand}.

%%K t02-17-1 p
17. --- Selline on inimkeha peaaegu täielik kirjeldus.

%%K t02-17-2 p
\textit{Märkus.} --- \textit{Väljendi «mul valutab... (pea jne.)» abil saab õpetaja
kogu seda sõnavara uuesti kasutada hea või halva} \VUgras{kehalise seisundi}
\textit{seisukohalt.}

%%K t02-tit-4 sub
\VUsaut{5.0mm}
\VUpk{10.2pt}{40}{Võimlemine.}
\VUsaut{2.4mm}
\VUcentre{10.2pt}{0}{\textit{(Ühe õpilase jutustus.)}}
\VUsaut{3.0mm}

%%K t02-18-1 p
18. --- Et olla painduvad, väledad ja terved, peame me harjutama oma jalgu ja käsi.
Seepärast me mängime ja teeme \VUgras{võimlemist}.

%%K t02-19-1 p
19. --- Nädala sees toimuvad need kaks harjutust gümnaasiumi õues (vaata tabelit
Nr. 1). Kuid neljapäeval ja pühapäeval läheme me suurele murule, kus meil on ruumi
joosta ja lõbutseda.

%%K t02-20-1 p
20. --- Meie õpetajad ja vanemad saadavad meid mõnikord sinna; kuid harjutusi juhib
ikkagi \VUgras{võimlemisõpetaja}\textsuperscript{(12)}.

%%K t02-21-1 p
21. --- Murul, sissekäigust vasakul, on \VUgras{võimlemisriistad}\textsuperscript{(1)}:
me ronime mööda \VUgras{redelit}\textsuperscript{(2)}, me teeme tirelit
\VUgras{rõngastel}\textsuperscript{(3)} ja \VUgras{trapetsil}\textsuperscript{(4)}, me
tõmbame end kätega üles \VUgras{köisredeli}\textsuperscript{(5)} või
\VUgras{sõlmedega köie}\textsuperscript{(6)} otsa.

%%K t02-22-1 p
22. --- Vahepeal hüppavad meie kaaslased üle \VUgras{puuhobuse}\textsuperscript{(7)}
või \VUgras{rööbaspuude}\textsuperscript{(11)}. Teised
\VUgras{poksivad}\textsuperscript{(8)} või \VUgras{vehklevad}\textsuperscript{(9)} maski
ja \VUgras{floretiga}. Lõpuks \VUgras{tõstavad} teised
\VUgras{kange}\textsuperscript{(10)}.

%%K t02-tit-5 sub
\VUsaut{4.6mm}
\VUpk{10.2pt}{40}{Mängud.}
\VUsaut{3.0mm}

%%K t02-23-1 p
23. --- Võimlejate ümber mängivad poisid ja tüdrukud mitmesuguseid
\VUgras{mänge}. Tüdrukud lõbustavad end oma
\VUgras{vitsarõngaga}\textsuperscript{(14)}; nad hüppavad
\VUgras{hüppenööriga}\textsuperscript{(13)} või kutsuvad oma vennad ja nõod
\VUgras{kroketi}\textsuperscript{(15)}-partiile. Neile valmistab rõõmu lüüa vastase
\VUgras{kuul} oma \VUgras{puuhaamriga} eemale just siis, kui too arvab kergesti
väravakesest läbi minevat!

%%K t02-24-1 p
24. --- Poisid eelistavad jõulisemaid mänge. Meelsasti mängivad nad
\VUgras{kurikaid}\textsuperscript{(16)}, mille nad osavalt \VUgras{kuuliga}\textsuperscript{(17)}
maha löövad. Nad ei põlga ka \VUgras{vurri}\textsuperscript{(18)} ja
\VUgras{bilbokeed}\textsuperscript{(19)} või isegi
\VUgras{konnamängu}\textsuperscript{(20)}. Kuid nende lemmikmängud on ikkagi
\VUgras{kuulikesed}\textsuperscript{(21)} ja \VUgras{seinapall}\textsuperscript{(22)},
sest \VUgras{kasvuhoone}\textsuperscript{(26)} sein sobib väga hästi kerget palli
tagasi põrgatama.

%%K t02-25-1 p
25. --- Väike Jaan, kes käib oma \VUgras{karkudel}\textsuperscript{(24)}, on väga osav
ja oskab väga hästi tasakaalu hoida. Tema lähedal mängivad mõned kaaslased
\VUgras{peitust}\textsuperscript{(25)} selles kasvuhoones ja \VUgras{heki} taga. Teised
eelistavad \VUgras{löögiäraarvamist}\textsuperscript{(28)} või
\VUgras{sulgpalli}\textsuperscript{(29)}, mis lendab ühelt \VUgras{reketilt} teisele.

%%K t02-noto-1 noto
\VUnotes{143.2mm}{%
(*) Laste- ja poistemängudel on sageli puhtrahvuslikud nimed. Me oleme püüdnud neid
idos nimetada nii hästi kui võimalik, kas lastes end juhtida tegevusest endast, mis
neid iseloomustab, või valides rahvusliku nime ühes või teises maas, kui see ei ole
liiga ebatäpne. Näiteks: \VUgras{vaadimäng} (saksa ja prantsuse) on
\VUgras{konnamäng} \textsuperscript{(20)} (inglise), milles mängijad püüavad heita oma
ümmargusi ketikesi läbi metallkonna suu, mis seisab suu avali augulisel kastil
(vaadil). \VUgras{Õlahüpe}\textsuperscript{(30)} erineb \VUgras{seljahüppest} ainult
selle poolest: pea kohalt hüppamiseks toetavad mängijad käed kaaslase õlgadele,
sellal kui «\VUgras{seljahüppes}» toetavad nad käed ainult tema seljale. --- Või siis
on mõned osavõtjatest kummargil, sellal kui teised hüppavad üle nende selja, toetades
käed õlale; esimesed peavad taluma hoopi kõverdumata, teised peavad hüppama tasakaalu
kaotamata (vaata tabelit ennast).%
}

%%K t02-26-1 p
26. --- Muru keskel on kaks puud. Ühe neist jalamil on seitse või kaheksa poissi
korraldanud \VUgras{õlahüppe}\textsuperscript{(30)}-partii (*). Seda mängu ei tunta
kõigis maades; kuid kõik teavad, mis on \VUgras{seljahüpe}, mida poisid seekord ei
mängi.

%%K t02-tit-6 sub
\VUsaut{5.0mm}
\VUpk{10.2pt}{40}{Mängud vabas õhus.}
\VUsaut{3.4mm}

%%K t02-27-1 p
27. \VUgras{Pimesikk}\textsuperscript{(31)} on samuti väga lõbus; tüdrukud ja poisid
panevad kordamööda \VUgras{silmasideme} silmadele ja püüavad kinni need osavusetud,
kes end tabada lasevad.

%%K t02-28-1 p
28. --- Muru keskel, siin on üks väga prantsuspärane, kuid pisut vägivaldne mäng:
see on \VUgras{karumäng}\textsuperscript{(32)}, mis on palju lärmakam kui nii vaikne
\VUgras{tuulelohe}\textsuperscript{(33)} mäng või kui inglise mäng
\VUgras{tennis}\textsuperscript{(34)}.

%%K t02-29-1 p
29. --- See viimane sport on suurte õpilaste rõõm. Meie õpetajad ise mängivad seda
meelsasti. See nõuab suurt osavust ja väledust ning mulle meeldib see väga.
\VUgras{Kiige}\textsuperscript{(35)} mängud jätan ma tüdrukutele.

%%K t02-30-1 p
30. --- Mulle ei meeldi üks teine inglise mäng, mis võib ohtlikuks muutuda.

%%K t02-30-1b p
Ma mõtlen \VUgras{jalgpalli}\textsuperscript{(36)}, mis rõõmustab mu vanemat venda. Ma
eelistan meie vana \VUgras{jooksumängu}\textsuperscript{(38)}, mis on niisama tervislik
ja tõesti vähem toores.

%%K t02-tit-7 sub
\VUsaut{4.6mm}
\VUpk{10.2pt}{40}{Rattasõit ja pallimängud.}
\VUsaut{3.0mm}

%%K t02-31-1 p
31. --- Samuti sõidan ma meelsasti \VUgras{jalgrattaga}\textsuperscript{(37)} ümber
muru.

%%K t02-32-1 p
32. --- Järgmisel aastal õpin ma \VUgras{ratsutama}\textsuperscript{(39)}. Kui ilus on
hobune, kes graatsiliselt sammub või traavib ja kes galopis üle takistuste hüppab!

%%K t02-33-1 p
33. --- Suvel õpime me \VUgras{ujumiskunsti}\textsuperscript{(40)}. Me sukeldume ja
supleme mõnuga jõe selges ja jahedas vees. Mõnikord laseme lõpuks lendu paberist
\VUgras{õhupalli}\textsuperscript{(41)}, mis tõuseb uhkelt õhku niipea, kui see on
sooja õhku täis.

%%K t02-34-1 p
34. --- On veel palju teisi mänge, kuid ma ei jõuaks neid kõiki üles lugeda, ja selle
lühikese kokkuvõtte järgi on näha, et neljapäeviti ei ole meie ilusal murul sugugi
igav.

%%K t02-34-2 p
N.-B. --- \textit{Õpetaja täiendab seda peatükki kergesti, kirjeldades
üksikasjalikumalt igat tähtsamat mängu.}
\VUsaut{6.0mm}
\VUfilet{22mm}
